% Einheit 1 von 4 – Abstand zweier Punkte (bank/abstaende/e1.jsonl, zusammenbau v0.1)
%% TODO Zeitmarke Einheit 1: Marken-Zeile nicht lesbar
%% TODO Zweigzeile Teil 1: was hier gelernt wird (Satz in Schülersprache aus dem Einheitstitel) – Einheit 1
\einheitenkopf[e1]{Einheit 1 von 4 · Abstand zweier Punkte}
\zweigzeile{Abitur GK}
\setcounter{aufgabe}{8}

%% TODO Titel: Ich-kann-Satz fehlt in der Bank (Platzhalter: Kettenname) – Nr. 9
%% TODO Anweisung: ein Satz über der ersten Teilaufgabe fehlt in der Bank – Nr. 9
\begin{aufgabe}{Hin oder zurück?}
\begin{teile}
\teil Hin oder zurück? Wie weit liegen A(2 | 0 | 1) und B(5 | 4 | 1) auseinander? Kreuze an, rechne nicht. \\ \kreuz{hin: einen Abstand berechnen} \\ \kreuz{zurück: aus einem Abstand einen Punkt oder Parameter bestimmen}
\end{teile}
\end{aufgabe}

%% TODO Titel: Ich-kann-Satz fehlt in der Bank (Platzhalter: Kettenname) – Nr. 10
%% TODO Anweisung: ein Satz über der ersten Teilaufgabe fehlt in der Bank – Nr. 10
\begin{aufgabe}{Abstand zweier Punkte}
%% TODO \rechenplatz{n}: Schrittzahl des Grundfalls fehlt in der Bank (nur bei fester Schreibform, 2.3 b)
\begin{teile}
\teil Abstand wovon zu was? Wie weit ist die Mastspitze S(3 | 2 | 7) vom Punkt T(0 | 6 | 7) entfernt? Kreuze das Verfahren an. \\ \kreuz{Punkt–Punkt: Betrag des Verbindungsvektors} \\ \kreuz{Punkt–Ebene: Hessesche Normalform} \\ \kreuz{Punkt–Gerade: Lotfußpunkt}
\teil A(2 | 1 | 3), B(4 | 3 | 4) – wie weit voneinander entfernt? d = \leerfeld
\teil Eine Drohne fliegt geradlinig in 4 s von P(5 | 20 | 30) nach Q(65 | 100 | 30); 1 LE = 1 m. Wie schnell ist sie in km/h? \hfill (Abitur 2018 GK) \\ v = \leerfeld km/h
\teil Ein Zylinder steht auf der xy-Ebene; sein Grundkreis hat den Mittelpunkt M(2 | 3 | 0) und den Radius 5, die Höhe ist 8. P(5 | 7 | 0) liegt auf dem Rand des Grundkreises. Welcher Punkt S des Deckkreisrands liegt P am nächsten, welcher Punkt T am fernsten? \hfill (Abitur 2020 LK) \\ S( \leerfeld | \leerfeld | \leerfeld ), T( \leerfeld | \leerfeld | \leerfeld )
\end{teile}
\begin{gleichungsraster}[2]
\gl{\text{Das Dreieck ABC mit A(0 | 0 | 0), B(2 | 1 | 2) und C(t | −2t | 0), t positiv, ist rechtwinklig in A und gleichschenklig. Welcher Wert von t? (Abitur 2026 GK)}} &
\end{gleichungsraster}
\begin{teile}
\teil Die Dreiecke $AB_sC_t$ mit A(0 | 0 | 0), $B_s(2s | 4s | 4s)$ und $C_t(2t | t | -2t)$, s und t positiv, sind rechtwinklig in A. Die Abbildung zeigt eine Gerade im s-t-Koordinatensystem. Weise nach: Die Punkte (s | t), für die das Dreieck gleichschenklig ist, liegen auf dieser Geraden. \hfill (Abitur 2026 LK)

\begin{ksys}[xmin=0,xmax=4,ymin=0,ymax=8,xlabel=s,ylabel=t,ablesen]
\gerade{2}{0}{}
\end{ksys}
\teil Der Wasserstrahl eines Rasensprengers folgt den Punkten $(3a | 4a | 1 + 2a - a^2)$, $a \ge 0$; 1 LE = 1 m, die xy-Ebene ist der Rasen, die Düse ist der Punkt für $a = 0$. Wie weit reicht der Strahl, gemessen auf dem Rasen vom Fußpunkt der Düse aus? Runde auf Zehntel Meter. \hfill (Abitur 2026 LK) \\ \leerfeld[m]
\teil Die Spitze eines Kamerakrans bewegt sich auf einem Kreis um H(3 | 1 | 4) mit Radius 5; der Kreis liegt in der Ebene L: $y + z = 5$; 1 LE = 1 m. Beurteile die Aussage: „Die Spitze erreicht eine Position mit der y-Koordinate 5.“ \hfill (Abitur 2026 LK)
\end{teile}
\end{aufgabe}

%% TODO Titel: Ich-kann-Satz fehlt in der Bank (Platzhalter: Kettenname) – Nr. 11
%% TODO Anweisung: ein Satz über der ersten Teilaufgabe fehlt in der Bank – Nr. 11
\begin{aufgabe}{Abstand zweier Punkte – Fehler finden, Begründen, Anwendung}
\begin{teile}
\teil Finde den Fehler. A(3 | 1 | −1), B(4 | 3 | 1). \rechnung{\overrightarrow{AB} &= (7 | 4 | 0) \\ |\overrightarrow{AB}| &= \sqrt{65} \approx 8{,}06}
\teil Warum ist die Formel für den Abstand zweier Punkte der Satz des Pythagoras im Raum? Begründe.
\teil Eine Seilbahn führt geradlinig von der Talstation T(0 | 0 | 450) zur Bergstation B(1\,200 | 900 | 1\,650); 1 LE = 1 m. Die Kabine fährt mit 5 m/s. Wie lange dauert eine Fahrt? \leerfeld[min]
\end{teile}
\end{aufgabe}
